\chapter{Trapped-Ion Quantum Computing}
\label{ch:theory}
\minitoc


\section{The Ion}
\label{sec:Transitions}
    \begin{figure}[ht]
        \begin{center}
        \noindent\includegraphics[width=\linewidth]{figures/pdf_figure/ion.pdf}
        \end{center}
        \caption{
        Electronic energy levels of \textsuperscript{40}Ca\textsuperscript{+},
        which will be used in this thesis. The levels are
        split by the Zeeman effect due to a $5.4$~G external magnetic field (they are shown explicitly only for the qubit manifolds). The
        transitions marked are required for cooling and control over the
        ion. The chosen qubit levels are labelled with $\ket{\downarrow}$, $\ket{\uparrow}$.
        }

    \label{fig:ion}
    \end{figure}

    \begin{figure}
        \vspace*{-0.5cm}
        \begin{center}
        \noindent\includegraphics[width=0.65\linewidth]{
            figures/pdf_figure/qp_gamma.pdf
            }
        \vspace*{-0.5cm}
        \noindent\includegraphics[width=0.65\linewidth]{
            figures/pdf_figure/qp_transition_spectrum_0.50.pdf
            }
        \vspace*{-0.5cm}
        \noindent\includegraphics[width=0.65\linewidth]{
            figures/pdf_figure/qp_transition_spectrum_0.25.pdf
            }
        \end{center}
        \vspace*{-0.5cm}
        \caption{
            Normalised quadrupole Rabi frequencies, $\Omega$, and frequency splittings for the $4S_{1/2} \leftrightarrow 3D_{5/2}$ Zeeman sublevels in \ca. Solid bars correspond to transitions with lower state $\ket{4S_{1/2},~m_j = -1/2}$, while open bars correspond to transitions with lower state $\ket{4S_{1/2},~m_j = +1/2}$.
            \textbf{a)} Relative Rabi frequencies as angle $\gamma$ between the polarisation vector and the magnetic field is varied. The k-vector of the light is assumed to be perpendicular to the B-field.
            \textbf{b)} Relative frequency splittings and strengths when $\gamma = \pi/2$. This is the polarisation chosen for the current and future experiments to reduce the number of transitions which can cause off-resonant interactions when the qubit defined by $\ket{4S_{1/2},~m_j = -1/2} \leftrightarrow \ket{3D_{5/2},~m_j = -5/2}$ is used.
            \textbf{c)} Relative frequency splittings and strengths when $\gamma = \pi/4$. At this angle, all transitions are accessible apart from those with $\Delta m_j = 0$. This is the polarisation chosen for our initial experiments to allow access to all transitions for characterisation.
            }
        \label{fig:quadrupole}
    \end{figure}

    


\section{Trapping Charged Particles}
\label{sec:trap_charge}


\section{Coherent Light-Matter Interactions}

\subsection{Rabi Flopping}
\label{sec:Rabi Flopping}

\subsection{Coupling to the Motion}
\label{sec:Motional Coupling}

\subsection{The Spin-Dependent Force}
\label{sec:SDF}

\subsection{Geometric Phase Gates}
\label{sec:Geometric Phase Gates}